% ============================================================
% §5.2 Q2结果分析 — MILP-MCLP-ES选址-规模优化 (v2: seed=42确定解)
% 竞赛: 电工杯 B题 (2026)
% ============================================================

\section{Q2: 选址与规模优化结果}

\subsection{最优选址-规模方案}

基于§3.2的MILP-MCLP-ES模型，以PuLP+CBC在120秒内收敛至gap$<1\%$的确定解（randomSeed=42），得120万元预算下F(大型)+H(中型)+J(中型)三站、总成本109万元、\textbf{10/10全覆盖}的最优方案（图\ref{fig:q2_network}），目标值54.225。站点运营指标见表\ref{tab:q2_stations}。

\begin{figure}[htbp]
\centering
\includegraphics[width=0.68\textwidth]{figures/figure_q2_network.png}
\caption{Q2最优选址-规模方案的空间网络图。蓝色节点为大型站(F)，橙色节点为中型站(H, J)，红色箭头表示服务辐射关系。}
\label{fig:q2_network}
\end{figure}

\begin{table}[H]
\centering
\small
\caption{Q2最优方案站点运营指标}
\label{tab:q2_stations}
\begin{tabular}{lcccccc}
\toprule
\textbf{站点} & \textbf{规模} & \textbf{建设成本(万元)} & \textbf{日容量(人次)} & \textbf{覆盖小区数} & \textbf{年利润(万元)} & \textbf{利润率(\%)} \\
\midrule
F & 大型 & 45 & 3000 & 4 & 497.6 & 26.1 \\
H & 中型 & 32 & 2000 & 3 & 306.7 & 24.7 \\
J & 中型 & 32 & 2000 & 3 & 293.4 & 24.4 \\
\midrule
\textbf{合计} & -- & \textbf{109} & \textbf{7000} & \textbf{10} & \textbf{1,097.7} & \textbf{25.2(加权)} \\
\bottomrule
\end{tabular}
\end{table}

\subsection{服务分配与满意度解构}

表\ref{tab:q2_assignments}列出最优分配与满意度分解。本方案以F+H+J一大型两中型（109万）替代C+J双大型（90万），追加H站实现对B、E的覆盖，代价为三站$S_2$因利用率超95\%全员触底0.600——以容量压榨换取地理全覆盖。

\begin{table}[htbp]
\centering
\caption{最优服务分配与满意度分解（全覆盖方案）}
\label{tab:q2_assignments}
\begin{tabular}{lcccccl}
\toprule
\textbf{需求小区} & \textbf{服务站} & \textbf{距离(m)} & $S_1$(距离) & $S_2$(响应) & $S$(综合) & \textbf{说明} \\
\midrule
A & J & 500 & 0.900 & 0.600 & 0.860 & 中距离，J站辐射 \\
B & H & 400 & 0.900 & 0.600 & 0.860 & 中距离，H站辐射 \\
C & F & 900 & 0.600 & 0.600 & 0.800 & 远距离触底，F站辐射 \\
D & J & 400 & 0.900 & 0.600 & 0.860 & 中距离，J站辐射 \\
E & H & 1000 & 0.600 & 0.600 & 0.800 & 边界距离，H站辐射 \\
F & F & 0 & 1.000 & 0.600 & 0.880 & 本小区自服务 \\
G & F & 500 & 0.750 & 0.600 & 0.830 & 中距离，F站辐射 \\
H & H & 0 & 1.000 & 0.600 & 0.880 & 本小区自服务 \\
I & F & 700 & 0.600 & 0.600 & 0.800 & 远距离，F站辐射 \\
J & J & 0 & 1.000 & 0.600 & 0.880 & 本小区自服务 \\
\bottomrule
\end{tabular}
\end{table}

\begin{figure}[H]
\centering
\includegraphics[width=0.92\textwidth]{figures/figure_q2_station_load.png}
\caption{各站点服务负荷与利用率对比。三站利用率均超过91\%，F站(95.9\%)已接近容量上限，120万元预算下不存在闲置容量。}
\label{fig:q2_load}
\end{figure}

\subsection{道法术三维解读}

Q2目标函数覆盖权重5.0与满意度权重0.5的比例为10:1，本质是对"不落一户"政策偏好的数学编码。CBC在39秒找到10/10覆盖可行解，120秒以Optimal状态退出，randomSeed=42确定的启发式分支顺序是避免局部最优的关键。最优方案呈三足鼎立结构：F大型站覆盖C-F-G-I四小区，H中型站覆盖B-E-H三小区，J中型站覆盖A-D-J三小区，三站日负荷各占容量上限的95.9\%、94.5\%、92.5\%（图\ref{fig:q2_load}），120万元预算下不存在闲置容量。确定解较非确定解（B小+C大+J大，8/10覆盖，目标值43.644）提升24\%，揭示MILP在非凸满意度景观上的多峰性质——多起点随机重启是回避局部最优的低成本策略\cite{liu2023}。Q2输出$y_{ik}^*$与$x_{ij}^*$作为Q3定价的固定输入，遵循"先选址、后定价"两阶段分离（合理性论证见§3.3）。
